\cotefile{13}
\scanpage{14}
Algiers Nov.~1, 1965\\
Dear Mumford,\\

I have become aware that I have been a bit rash with the conjectures I told you about Hodge structures, as in the form I stated them they contradict \unsure{Tates|Tate's} conjecture, \deleted{\unsure{w|a}}which implies the following: if $X$ is a smooth, connected \uline{simply connected} scheme over $\mathbb{C}$, and $\unsure{\mathscr{V}|V}$ a polarized complex analytic family of Hodge structures parametrized by $X^{\mathrm{hol}}$, then $\mathscr{V}$ is ``algebraic'' only if it is a \uline{constant} family. Now it is easy to get examples of $\mathscr{V}$'s, (with $X$ say \uline{any}\unsure{.|{}} \unsure{homogeneous|homogenous} space under an affine group over $\unsure{\mathbb{C}|C}$)
% archive sheet number: 13
\scanpage{15}
which are \uline{non} constant, therefore (one hopes!) not ``algebraic'' (as one would like Tate's \unsure{conjecture|conjectures} to hold).

However, one should try to test if one really can\uline{not} get an algebraic $\unsure{\mathscr{V}|V}$ that way. To start, I wonder if a single Hodge-structure $V$ with \uline{even partial degrees} (therefore giving rise to a ``Hodge-group'' $G$ s.th. $G_{\mathbb{R}}$ is \uline{compact}) can be algebraic \inserted{(except when $G$ is \unsure{comm.|comm})}.

For instance, do you know an algebraic smooth projective surface $Y$ over $\unsure{\mathbb{C}|C}$, such that ${\unsure{H|H\!V}}^{11}(Y,\mathbb{C})$ is spanned by algebraic cycles, and $\dim H^{0,2}(Y,\mathbb{C})\geq 2$ \unsure{?!|??} I would appreciate your comments! \quad Yours\unsure{.|{}}\\
A Grothendieck\\

P.S. What about your $\Theta$-\unsure{fns|funs}, are you going to write a book?
% archive sheet number: 14
\scanpage{16}
P.S. Here what are the modular varieties one gets from considerations of motive-theoretic galois theory.

\[
G\deleted{_{\unsure{\mathbb{R}|\mathbb{Q}}}}
= \text{reductive connected alg. group defined over }\mathbb{Q}
\]

\[
\begin{aligned}
U_{\mathbb{R}}={}&\text{dimension one twisted real Torus, with }U_{\mathbb{R}}(\mathbb{R})\\
&\text{identified with the group of complex nbs of}\\
&\text{module }1
\end{aligned}
\]

\[
\begin{aligned}
i_0:U_{\mathbb{R}}\longrightarrow G_{\mathbb{R}}\quad &\text{homomorphism of algebraic}\\
&\text{groups over }\mathbb{R},\ \text{s.th. }i_0(-1)\text{ is}\\
&\text{central, and s.th. \deleted{\unsure{t|a}}the \unsure{centralisator|centralizator}}\\
&K_{\mathbb{R}}\text{ of }i_0(\sqrt{-1})\text{ is such that}\\
&K_{\mathbb{R}}(\mathbb{R})^{\unsure{0|\circ}}\text{ is \uline{maximal} \uline{compact} in }G_{\mathbb{R}}(\mathbb{R})^{0}.
\end{aligned}
\]

Moreover, for every $\mathbb{R}$-simple component of $G_{\mathbb{R}}$, the corresponding component of $i_0$ is not trivial.

\[
C_{\mathbb{R}}=\text{\uline{Centr}}_{G_{\mathbb{R}}}(i_0)\subset K_{\mathbb{R}}
\]

\[
\deleted{\unsure{\gamma|\mathscr{T}}}Q=G_{\mathbb{R}}/C_{\mathbb{R}}
\]

\[
\text{``\unsure{Sigel|Siegel}-space''}=\gamma=G(\mathbb{R})^{0}/G(\mathbb{R})^{0}\cap C_{\mathbb{R}}(\mathbb{R})
\overset{\substack{\inserted{\text{(a connected component)}}\\ \deleted{\text{\unsure{open and closed}}}}}{\subset}Q(\mathbb{R})
\]

N.B. $\gamma$ has a natural complex structure (inherited from \inserted{one on} $Q(\mathbb{R})$ in fact) invariant by operations of $G(\mathbb{R})$. Every linear representation (over $\mathbb{Q}$) of $G$ in\deleted{to} a vector space $V$ (over $\mathbb{Q}$), together with a lattice $V_0$ of $V$, defines a \unsure{``|‘}\uline{complex} \uline{analytic} family \inserted{$\unsure{\mathscr{V}|V}$} of Hodge structures\textsuperscript{\ensuremath{\delta}}\margin{$\delta$\;(\deleted{\unsure{perhaps even|perhaps}} \inserted{necessarily} globally polarizable -- to get a polarization of the family, \unsure{one} picks a ``form of polarization'' $\unsure{\varphi|\phi}$ \uline{on $V$} \deleted{\unsure{,|;}} invariant by $G$).} on $\gamma$, at least when one fixes a ``total degree'' $n$
% archive sheet number: 15
\scanpage{17}
and assumes that $i_0(-1)$ operates on $V$ as \inserted{homothety} $(-1)^n$. (It is not hard to describe $\unsure{\gamma|\mathscr{T}}$ by universal propert\deleted{\unsure{ies|y}}y as \deleted{\unsure{being}} representing some functor \inserted{$F(X)$, $X$ \deleted{\unsure{al|an}}complex analytic,} expressing that every $V$ \inserted{as above} should define, in an additive and multiplicative way for varying $V$, a Hodge-structure \inserted{$\unsure{\mathscr{V}|V}$} over $X$). If $\Gamma\subset G(\deleted{\unsure{\mathbb{Q}|\mathbb{R}}}\unsure{\mathbb{R}|\mathbb{Q}})$ is the group leaving $V_0$ fixed, then $\Gamma$ operates on $(\mathscr{V},\gamma)$.
% margin note written sideways
\margin{Note that in general $\gamma$ is \uline{not} \unsure{riemanian|riemannian} symm., but is a fiber space over a R.S. space, with fibers which are projective complex \unsure{homogeneous|homogenous} spaces under $K_{\mathbb{C}}$. This implies at least that $\gamma$ is \uline{simply} \uline{connected}, hence is the \unsure{univ.|univ-} covering of $M=\gamma/\Gamma$ if $\Gamma$ operates freely}

The modular varieties I expect to carry algebraic structure\deleted{\unsure{s}}\unsure{{}|s} \uline{over} $\mathbb{Q}$ (and morally, $\mathbb{Z}$) are the varieties $\boxed{\gamma/\Gamma}$, where $\Gamma\subset G(\mathbb{Q})\cap G(\mathbb{R}^{\unsure{0|\circ}})$ \deleted{\unsure{is}} is commensurable with $G(\mathbb{Z})$ [for some matrix representation of $G$, \deleted{\unsure{as|so}} giving a meaning to $G(\mathbb{Z})$]. However, it is possible that one will have to impose on the data some extra conditions, either arithmetic (involving existence of ``Frobenius elements'' in $\deleted{\unsure{G(\mathbb{Q})|G}}\in\unsure{.|{}}\;I(\mathbb{Q})$, $I=G/\text{\unsure{élat.|élém.} \unsure{aut.|ant.}}$),\unsure{.|{}} or geometric \inserted{on $G_{\mathbb{R}}$} (such as: $G_{\mathbb{R}}$ has \uline{no compact factor}). A typical case which fits in the general \deleted{\unsure{picture|pictures}} \inserted{description}, but for which I have no evidence so far if $\gamma/\Gamma$ has
% archive sheet number: 16
\scanpage{18}
an algebro-geometric interpretation in terms of \inserted{moduli for} motives, is the following:

start with a vector space of finite \unsure{:|{}}\unsure{dimension-|dimension} $V$ over $\mathbb{Q}$, with \unsure{an|a} fundamental bilinear form $\unsure{\varphi|\phi}$ given, symmetric or alternating, take $G=\mathrm{SO}(\varphi)$, take any $i_0:U_{\mathbb{R}}\to G_{\mathbb{R}}$ s.th. $i_0(-1)=\mathrm{id}$ (if $\varphi$ symmetric), $i_0(-1)=-\unsure{\mathrm{Id}|\mathrm{id}}$ (if $\varphi$ \deleted{\unsure{a|o}} skew-symmetric), \deleted{$\unsure{i_0|\mathrm{id}}(\sqrt{-1})\unsure{=|\ne}\unsure{1|{-1}}$} \deleted{$i_0(i)$ \unsure{non} trivial.} \inserted{and a positivity condition of $\varphi$ relative to the bigrading of $V_{\mathbb{C}}$ defined by $i_0$, which I will not write down.}

Then $\unsure{\gamma|\mathscr{T}}$ is a direct generalization of Siegel's half-plane, and classifies all Hodge-structures \deleted{\unsure{under}} with underlying $\deleted{\unsure{V|Q}}\mathbb{Q}$-vector space $V$, \deleted{\unsure{(|a}} i.e. all bigradings of $V_{\mathbb{C}}$, \deleted{\unsure{having}} \inserted{satisfying} the usual symmetry condition for complex conjugation, which are compatible with $\varphi$ (\unsure{i.e|i.e.} $\varphi$ is a polarization), and for which the $\dim V_{\mathbb{C}}^{pq}$ have given values \unsure{.|{}} (depending on the weights of $U_{\deleted{\mathbb{R}}\mathbb{C}}$ operating on $V_{\mathbb{C}}\ \ldots$). The trouble is that, \deleted{\unsure{if|of}} as these special types of $\gamma$, $\gamma/\Gamma$ are in a way ``universal'' for all others
% archive sheet number: 17
\scanpage{19}
(namely the $\unsure{\mathscr{V}|V}$'s considered above are induced from one's on such special type of $\unsure{\gamma|\mathscr{T}}$), if the ``modular family'' $\mathscr{V}$ of Hodge structures on \inserted{such a} $\gamma/\Gamma$ was ``algebraic'' i.e. came from a relative motive over the complex analytic space $\gamma/\Gamma$, \inserted{(or even on some underlying structure of alg. var. over $\mathbb{C}$),} the same would be true for \uline{every} modular family on some $\gamma/\Gamma$.

Taking the case $G\deleted{_{\mathbb{R}}}(\mathbb{R})$ compact, \deleted{\unsure{this|this is}} and \unsure{thus|then} $\gamma/\Gamma=\gamma={}$complex projective \unsure{homogeneous|homogenous} space $X$ under $G_{\mathbb{C}}$ by Borel, we would get a \inserted{polarized} motive over the analytic space $\gamma$ corresponding to $X$. If we admit the \inserted{GAGA} yoga (coming from \inserted{polarized} abelian varieties)\unsure{.|{}} that such a family must also be a motive over $X$, \inserted{(or the above strengthening of the original assumption),} we \deleted{would} get a contradiction with Tate's \unsure{conjectures|conjecture} \deleted{\unsure{as|a}} (\deleted{\unsure{bec|as}}because $X$ is simply connected, as well known). So I really do not know what to believe!
% archive sheet number: 18
\scanpage{20}

To come back to the general case, maybe I should add that a \uline{\unsure{necessary|necesary|neccessary|neccesary}} condition for $\unsure{\gamma|\mathscr{T}}/\Gamma$ to be algebraic and the \unsure{Hodge-structures|Hodge-structure} $\unsure{\mathscr{V}|V}/\Gamma$ over $\gamma/\Gamma$ algebraic too, is (\uline{granting} Tate's and Hodge's \unsure{conjectures|conjecture}): for every $g\in G(\mathbb{R})^{\unsure{0|\circ}}$, and every $\Gamma_0$ of finite index in $G(\mathbb{Z})$, the smallest \inserted{alg.} subgroup $H$ of $G$ such that $H_{\mathbb{R}}$ contains $g[i_0(U_{\mathbb{R}})]g^{-1}$ \uline{and} $\unsure{\Gamma|T}_0$ is $G$ itself.

\deleted{I did \unsure{not} try to test directly this condition, even in the case of the \uline{usual} Siegel modular space; if it should be false in that case, as we know that $\gamma/\Gamma$ \uline{is} algebraic and that so are the $\mathscr{V}/\Gamma$ over it, it would follow that either Hodge's or Tate's conjecture is wrong\unsure{.|{}}}

\margin{Question: If $G_{\mathbb{R}}$ is semi-simple without compact factors, is $G(\mathbb{Z})$ always Zariski-dense in $G$\unsure{?!|??}}

By the way, did you make out if for the modular varieties $\gamma/\Gamma$ of your
\[
\left\{\begin{aligned}
&\text{Boulder talk (which I just read),}\\
&\text{the image of }\gamma/\Gamma\text{ in the usual Siegel}\\
&\text{modular variety is an algebraic variety,}\\
&\text{or at least constructible?}
\end{aligned}\right.
\]

% archive sheet number: 19
